%! Project: Design and Conduct a Survey — Unit 1, Lesson 1.8 — Student Packet Cover Sheet
% PILOT SHAPE (2026-09-06): modeled on AP Statistics lesson 1.4 minus its AP Practice page.
% THREE packet rows — Warm-Up, Guided Notes & Practice, Homework. The homework is the individual
% practice, started in class, and it is SCORED, so its cell is a \blank{}, never NA. This is
% the unit's project lesson: the Guided Practice row of the notes is the class's own survey,
% and the homework's last item is the project deliverable (the four-sentence report).
\documentclass[12pt]{article}
\usepackage{saar-article}
\usepackage{saar-boxes}
\usepackage{ltablex}
\keepXColumns
% Measured banner: the band grows to fit the title block, so a title long enough to wrap grows
% the band rather than running out of it (the stats course's \coverbanner; saar-article does not
% carry one, so it is defined here — every cover copies this block verbatim).
\newsavebox{\bannerbox}
\newlength{\bannerht}
\newlength{\coverbannerpad}
\setlength{\coverbannerpad}{0.16in}       % breathing room above and below the text
\newcommand{\coverbanner}[2]{%
  \sbox{\bannerbox}{%
    \begin{minipage}{\dimexpr\paperwidth-1.4in\relax}%
      \centering
      {\color{white}\textbf{\LARGE \CourseName}}\\[4pt]
      {\color{frostmid}\large\textbf{#1}}\\[6pt]
      {\color{white}\textbf{\large #2}}%
    \end{minipage}}%
  \setlength{\bannerht}{\dimexpr\ht\bannerbox+\dp\bannerbox+2\coverbannerpad\relax}%
  \noindent\begin{tikzpicture}[remember picture, overlay]
    \fill[cerulean] (current page.north west)
      rectangle ([yshift=-\bannerht]current page.north east);
    \node[anchor=north, inner sep=0pt]
      at ([yshift=-\coverbannerpad]current page.north) {\usebox{\bannerbox}};
  \end{tikzpicture}\par
  % The text body starts 0.6in below the page edge (geometry top=0.6in), so reserve only what
  % the band overruns that, plus a gap under the band.
  \vspace*{\dimexpr\bannerht-0.6in+0.16in\relax}%
}

\begin{document}

% Full-bleed cerulean banner, sized to the title block.
\coverbanner{Unit 1}{Lesson 1.8 \quad Project: Design and Conduct a Survey}

\vspace{0.06in}
% NAMESTRIP: this is the ONLY component that carries the name row.
\namedateperiod
\vspace{0.18in}

% GRADUAL RELEASE: the vocabulary is taught in the guided notes, so the targets name it
% plainly. The standard codes (AFDA.DA.2d, 2h, 2i, 2j; PS.DC.2d) stay in the teacher-facing plan.
\begin{learningtargetbox}
\small
\begin{enumerate}[leftmargin=1.5em, itemsep=5pt]
  \item \ldots{} plan a survey through the four stages of the \textbf{statistical cycle}
        and name what each stage hands in.
  \item \ldots{} write a \textbf{survey instrument} whose \textbf{closed-response items}
        are neutral, ask one thing, and offer choices that cover everyone without
        overlapping.
  \item \ldots{} turn a tally into \textbf{relative frequencies} and read them off a
        \textbf{bar graph}.
  \item \ldots{} report a conclusion about the people my design actually covers ---
        and explain why a \textbf{convenience sample} that happens to agree with a
        chance sample has proved nothing.
\end{enumerate}
\end{learningtargetbox}

\vspace{0.14in}

\begin{tocbox}
\small
\renewcommand{\arraystretch}{1.5}
\begin{tabularx}{\linewidth}{@{} c l X r @{}}
\rowcolor{frost}
\textbf{\#} & \textbf{Component} & \textbf{Description} & \textbf{Score} \\
\midrule
1 & Warm-Up & The Student Council's finished survey: a percent, a prediction, and the
             two groups hiding in the paragraph & \blank{1.2cm} \\
2 & Guided Notes \& Practice & The cycle as a checklist; four rules for an instrument;
             counts to percents to a bar graph; whom a report may describe ---
             then the survey this class runs on itself & \blank{1.2cm} \\
% Row 3 is the lesson's GRADED practice, started in class. Packet pages are the default; a
% Desmos activity may replace them on a given day, decided at Close & Assign. The row and its
% score cell never change.
3 & Homework & Millbrook Public Library: a phone survey and a front-desk clipboard ---
             then your class's own four-sentence report ---
             \textbf{scored, started in class, due the first class after two study halls} & \blank{1.2cm} \\
\midrule
  & \multicolumn{2}{r}{\textbf{Total}} & \blank{1.2cm} \\
\end{tabularx}
\end{tocbox}

\vspace{0.14in}

% Keep in Mind carries the lesson's CONTENT — the rule, the test, the trap — never its process.
\begin{remindbox}
\small
A \textbf{relative frequency} is a count divided by the total, and the four add to $100\%$
only when the \textbf{instrument} gave every person exactly one box to check. A sample percent
may be scaled to the population only when \textbf{chance} chose the sample --- ask \emph{who
chose these people?} before you multiply. \textbf{Watch out:} a class surveying itself is a
\textbf{convenience sample}, and landing near the council's number is luck, not evidence ---
the same class could have missed by twenty points and nobody would know. Name the limit in
your report; that is part of the answer, not an apology.
\end{remindbox}

\end{document}
